\documentclass[
  12pt,
  twoside,
  paper=a4,
  DIV=15,
  BCOR=12mm,
  abstract=true,
  headsepline,
  headings=normal
]{scrreprt}

% Compile with LuaLaTeX (recommended) or XeLaTeX.
% Keep this file stable; put chapter-specific macros into the chapter file.

\usepackage[onehalfspacing]{setspace}
\usepackage[T1]{fontenc}
\usepackage{lmodern}

% Language / typography
\usepackage[main=english,ngerman]{babel}
\usepackage{csquotes}
\usepackage{microtype}

% Graphics / tables
\usepackage{graphicx}
\graphicspath{{figures/}}
\usepackage{booktabs}
\usepackage{siunitx}

% Code (fast, no shell-escape)
\usepackage{listings}
\usepackage[dvipsnames]{xcolor}
\lstdefinestyle{code}{
  basicstyle=\ttfamily\small,
  columns=fullflexible,
  frame=single,
  breaklines=true,
  keywordstyle=\color{NavyBlue},
  commentstyle=\color{Gray},
  stringstyle=\color{OliveGreen},
}
\lstset{style=code}

% Algorithms (optional; remove if you don't use it)
\usepackage[noend]{algorithm2e}

% TODO notes (disable in final by switching to \usepackage[disable]{todonotes})
\usepackage[colorinlistoftodos,prependcaption,textsize=small]{todonotes}

% KOMA-friendly headers/footers
\usepackage[automark,headsepline]{scrlayer-scrpage}
\clearpairofpagestyles
\ihead*{\headmark}
\ofoot*{\pagemark}
\renewcommand*{\chapterpagestyle}{scrheadings}

% References
\usepackage[
  backend=biber,
  style=apa,
  sorting=ynt
]{biblatex}
\addbibresource{bib/references.bib}

% Links (load late)
\usepackage[colorlinks=true,urlcolor=NavyBlue,linkcolor=black,citecolor=BlueViolet]{hyperref}


\begin{document}

\begin{titlepage}
\centering

{\Large Master Thesis\par}
\vspace{1cm}
{\huge\bfseries ECG-Based Chagas Disease Detection with Deep Learning: Data Harmonization and Loss Function Strategies\par}
\vspace{1cm}

{\large Author: Florian Herzler\par}
{\large Supervisors: Dr. Sören Lukassen, Prof. Dr. Alexander Meyer\par}
\vfill

{\large \today\par}
\end{titlepage}


\includepdf[pages=1,fitpaper=true,pagecommand={\thispagestyle{empty}}]{figures/selbststaendigkeitserklaerung_fh.pdf}
\clearpage
\pagestyle{scrheadings}
\pagenumbering{roman}

\chapter*{Abstract}
\addcontentsline{toc}{chapter}{Abstract}
This thesis investigates ECG-based Chagas disease screening under severe class imbalance, label heterogeneity, and multi-source domain shift.
The work combines three datasets (CODE-15, SaMi-Trop, PTB-XL) and evaluates a fixed 1D ResNet-18 pipeline across three preprocessing regimes (\texttt{bp}, \texttt{bp-sc}, \texttt{bp-sc-norm}) and three learning tracks: end-to-end supervised classification (Track~1), supervised contrastive representation learning (Track~2), and linear probing of pretrained encoders (Track~3).
Evaluation follows a four-stage analysis workflow: (1) clinical ranking utility (primary endpoint: TPR@5\%), (2) embedding-health diagnostics, (3) cross-model ranking consistency and robust rank aggregation (RRA), and (4) physiological plausibility using ST-DFT-LRP.

Across fold~4 experiments, preprocessing was the dominant determinant of performance: \texttt{bp} and \texttt{bp-sc} consistently outperformed \texttt{bp-sc-norm}, while \texttt{bp-sc-norm} showed collapse-like embedding behavior and markedly weaker screening utility.
Embedding diagnostics supported these findings: class-separation metrics (especially CAC for class~1) aligned better with downstream TPR@5\% than uniformity-style metrics alone.
Cross-model agreement analyses showed structured clustering by preprocessing regime in global ranks, but substantially lower stability at the top-5\% decision boundary, motivating consensus-based sample selection.
RRA produced robust cross-model candidate sets that were informative for qualitative analysis but not identical to per-model score extremes.
Lead-level ST-DFT-LRP summaries highlighted consistent precordial emphasis and preprocessing-dependent attenuation of disease-contrast relevance patterns.
Perturbation checks provided supportive, though heterogeneous, faithfulness evidence for selected explanations.

Methodologically, the thesis shows that in this setting robust screening behavior is primarily controlled by data harmonization strength and objective design, not backbone changes. At the same time, results are constrained by single-fold evaluation, mixed label reliability, and model-pool dependence of consensus procedures. The findings provide a reproducible analysis framework and concrete hypotheses for future work on stronger augmentations, architecture transfer, and external validation.
\cleardoublepage

\tableofcontents

\pagenumbering{arabic}
\chapter{Introduction}

\section{Clinical Motivation}
Chagas disease is a neglected tropical disease caused by \textit{Trypanosoma cruzi} infection and can progress to chronic Chagas cardiomyopathy, which is associated with conduction abnormalities, arrhythmias, and heart failure.
Early identification of infected individuals enables confirmatory serological testing and timely follow-up, but population-wide screening is challenging in resource-constrained settings.

\section{Why the ECG?}
The 12-lead electrocardiogram (ECG) is inexpensive, widely available, and directly reflects cardiac electrical activity.
Since chronic Chagas disease can manifest with characteristic electrophysiological changes, the ECG is a plausible first-line screening modality to \emph{prioritize} who should receive confirmatory serology \parencite{jidling_screening_2023}.
Importantly, an ECG-based model is expected to be insensitive to early infection without cardiac involvement; in that setting the ECG may appear normal and serology remains the definitive test.

\section{Task Origin: PhysioNet/Computing in Cardiology Challenge 2025}
The experimental classification task in this thesis is derived from the George B.~Moody PhysioNet Challenge 2025, which focuses on detecting Chagas disease from the ECG \parencite{2025ChallengeCinC}.
We therefore interpret the task as learning a \emph{risk score for Chagas-related ECG manifestations} (as represented by the available labels), rather than a universal screen for all infected individuals irrespective of disease stage.
The challenge emphasizes performance at a low false-positive rate, motivating metric-aligned modeling choices.

\section{Why Deep Learning?}
ECG-based detection requires modeling subtle, distributed waveform patterns across leads and time.
Deep learning methods, particularly convolutional architectures, can learn such representations directly from raw waveforms and have shown strong performance for ECG interpretation tasks \parencite{jidling_screening_2023}.

\section{Key Challenges: Dataset Composition}
Two properties of the training data complicate optimization and interpretation:
\begin{itemize}
  \item \textbf{Class imbalance}: positives are rare, so standard objectives can be dominated by easy negatives.
  \item \textbf{Dataset bias and label noise}: the pooled dataset combines multiple sources with different acquisition pipelines and population characteristics, and some labels (e.g., in large-scale cohorts such as CODE-15\%) may be less reliable than clinically adjudicated endpoints \parencite{ribeiro_code-15_2021,jidling_screening_2023}.
\end{itemize}
These issues can lead to inflated apparent performance on in-distribution validation while degrading generalization across sources.

\section{Contributions and Thesis Structure}
This thesis studies practical strategies to improve robustness and embedding quality under the above constraints:
\begin{itemize}
  \item \textbf{Track 1}: supervised classification with metric-aligned loss functions and controlled augmentations.
  \item \textbf{Track 2}: supervised contrastive / prototype-based training to shape the embedding space and quantify representation quality.
  \item \textbf{Track 3}: downstream probing of pretrained encoders to connect embedding quality to classification performance.
\end{itemize}
Chapter~\ref{ch:methods-results} details the configuration system and experimental variables used in these tracks.

\chapter{Background}

This chapter introduces the core concepts needed to understand the experimental design in Chapter~\ref{ch:methods-results}.
The work is conducted in the context of the PhysioNet/Computing in Cardiology Challenge 2025 on Chagas disease detection from ECGs \parencite{2025ChallengePreprint,2025ChallengeCinC}, building on prior deep-learning approaches for ECG-based screening \parencite{jidling_screening_2023}.

\section{ECGs, Leads, and Preprocessing}
We work with standard 12-lead resting ECGs. The underlying ECG sources include large-scale cohorts such as CODE-15 \parencite{ribeiro_code-15_2021}, SaMi-Trop \parencite{ribeiro_sami-trop_2021}, and PTB-XL \parencite{wagner_ptb-xl_nodate}, distributed via PhysioNet \parencite{goldberger2000physiobank}.
The dataset is preprocessed into multiple variants to study the trade-off between signal standardization and potential loss of physiologically meaningful variation:
\begin{itemize}
  \item \textbf{bp}: band-pass filtered signals (approximately linear processing).
  \item \textbf{bp\_sc}: additional soft clipping (nonlinear amplitude transform).
  \item \textbf{bp\_sc\_norm}: additional z-score normalization (per-record standardization).
\end{itemize}
These variants are treated as different experimental regimes rather than interchangeable inputs.

\section{Learning Tracks}
The thesis is organized into three learning tracks:
\begin{enumerate}
  \item \textbf{Track 1 (classification)}: supervised training to predict the target label.
  \item \textbf{Track 2 (projection/embeddings)}: supervised contrastive learning to shape the embedding space.
  \item \textbf{Track 3 (classification with pretrained encoder)}: linear probing (or partial fine-tuning) on an encoder pretrained in Track~2.
\end{enumerate}

\section{Imbalanced Binary Classification and Loss Functions}
The label distribution is highly imbalanced. We therefore focus on losses that emphasize the minority class and/or ranking quality:
\begin{itemize}
  \item \textbf{Weighted BCE}: binary cross entropy with a positive class weight to counter imbalance.
  \item \textbf{Focal loss}: down-weights easy examples and focuses gradients on hard examples (controlled by $\gamma$) \parencite{focal_2018}.
  \item \textbf{Ranking-aware Tversky (RAT)}: a Tversky-style objective \parencite{tversky_2017} computed on a soft top-$k$ region of the batch (motivated by the challenge metric), with auxiliary regularizers for stability \parencite{rat2025}.
\end{itemize}

\section{Representation Learning: Supervised Contrastive Learning and Prototypes}
Track~2 uses a projection head and contrastive objectives to encourage embeddings of the same class to cluster and different classes to separate. We consider:
\begin{itemize}
  \item \textbf{Standard SupCon}: supervised contrastive loss over augmented views.
  \item \textbf{SupCon with imbalance-aware sampling/weighting}: optional modifications intended to reduce dominance of the majority class in binary imbalanced settings \parencite{ttc_2025}.
  \item \textbf{Prototype-based supervision}: class prototypes that act as anchors in representation space.
\end{itemize}

\section{Embedding Diagnostics and Early Stopping}
For Track~1, we monitor a task metric aligned with the challenge objective (TPR@5\% thresholding).
For Track~2, we rely on embedding-space diagnostics such as class alignment consistency (CAC) and qualitative visualization (e.g., UMAP) to judge whether representations improve.

\chapter{Methods}
\label{ch:methods}

\section{Configuration System}
Experiments are implemented in PyTorch \parencite{pytorch_2019} using PyTorch Lightning \parencite{pytorch_lightning_2025}.
Configuration is managed through composable YAML files, separating shared defaults (backbone, optimizer, data loading) from track-specific settings (augmentations, loss functions).
An individual run combines a base configuration, a track configuration, and an optional loss configuration, with specific parameters overridden via command line arguments.

\section{Model Architecture}
A 1D adaptation of the ResNet-18 architecture \parencite{resnet_2015} that operates directly on 12-lead ECG waveforms (Figure~\ref{fig:model_arch}) was implemented.
The network consists of an initial convolutional layer followed by four residual stages, each containing two blocks with 1D convolutions and internal temporal downsampling.
Global average pooling aggregates the features into a fixed-size vector.
Depending on the track, this backbone is followed by either a classification head (for classification learning) or a projection head (for contrastive learning).

\begin{figure}[htbp]
  \centering
  \resizebox{0.6\textwidth}{!}{%
    \begin{tikzpicture}[node distance=4mm,>=latex]
      \tikzset{
        arrow/.style={-Latex, thick},
        annot/.style={font=\footnotesize},
        op/.style={draw, rounded corners=2pt, thick, align=center, minimum height=5mm, minimum width=18mm, fill=gray!10},
        stem/.style={draw, rounded corners=2pt, thick, align=center, minimum height=5mm, minimum width=22mm, fill=gray!10},
        blk/.style={draw, rounded corners=2pt, thick, align=center, minimum height=5mm, minimum width=18mm},
        head/.style={draw, rounded corners=2pt, thick, align=center, minimum height=8mm, minimum width=18mm},
        tiny/.style={draw, rounded corners=2pt, thick, align=center, minimum height=4.5mm, minimum width=16mm, font=\footnotesize},
      }
      % =========================================================
      % Overview
      % =========================================================
      \node[op] (x) {12-lead ECG};
      \node[stem,below=3mm of x] (conv) {Conv-BN-ReLU\\MaxPool};
      \node[blk,below=3mm of conv,fill=orange!20] (resblk) {ResBlk $\times 4$};
      \node[op,below=3mm of resblk] (drop) {GAP-Flatten\\Dropout};

      \draw[arrow] (x) -- (conv);
      \draw[arrow] (conv) -- (resblk);
      \draw[arrow] (resblk) -- (drop);

      % --- split to heads without overlap ---
      \node[head,below=7mm of drop, xshift=-12mm, anchor=north,fill=blue!30] (proj) {Projection\\Head};
      \node[head,below=7mm of drop, xshift= 13mm, anchor=north,fill=blue!30] (clf)  {Classification\\Head};
      \draw[arrow] (drop) -- (proj);
      \draw[arrow] (drop) -- (clf);

      % =========================================================
      % Residual block (zoom)
      % =========================================================
      \node[draw, rounded corners=3pt, thick, inner sep=5pt, fill=orange!20,
        right=8mm of resblk, yshift=+3mm,
        minimum width=27mm, minimum height=57mm] (inset) {};
      \node[annot,font=\bfseries,anchor=north] (t_res)
      at ($(inset.north)+(0,-0.2mm)$) {ResBlk};

      \node[tiny,fill=blue!10,below=8mm of inset.north] (c1) {Conv $3{\times}1$};
      \node[tiny,fill=blue!10,below=3mm of c1] (bn1) {BN};
      \node[tiny,fill=red!20,below=3mm of bn1] (r1) {ReLU};
      \node[tiny,fill=blue!10,below=3mm of r1] (c2) {Conv $3{\times}1$};
      \node[tiny,fill=blue!10,below=3mm of c2] (bn2) {BN};
      \node[tiny,fill=red!20,below=3mm of bn2] (rout) {ReLU};

      \coordinate (inflow) at ($(c1.north)+(0,8mm)$);
      \draw[arrow] (inflow) -- (c1.north);
      \draw[arrow] (c1) -- (bn1);
      \draw[arrow] (bn1) -- (r1);
      \draw[arrow] (r1) -- (c2);
      \draw[arrow] (c2) -- (bn2);
      \draw[arrow] (bn2) -- (rout);

      % skip connection
      \draw[very thick,-Latex]
      ($(inflow)+(0,-1.5mm)$) |- ($(inflow)+(-12mm,-1.5mm)$) |- ($(rout.west)+(-1.5mm,0)$) |- (rout.west);
    \end{tikzpicture}
  }%
  \caption[Model architecture]{Model architecture with classification and projection heads (left) and the residual block definition (right).}
  \label{fig:model_arch}
\end{figure}

To limit the experimental scope, the architecture is fixed across all runs; batch normalization is used. This ensures that performance differences can be attributed to preprocessing, augmentations, and objective functions rather than model capacity.

\section{Data Regimes}
\label{sec:data_regimes}
We investigate three preprocessing levels to isolate the effect of harmonization:
\begin{enumerate}
  \item \textbf{bp}: Band-pass filtered signals ($0.67$--$45$\,Hz).
  \item \textbf{bp\_sc}: Band-pass filtering plus soft clipping to handle outliers.
  \item \textbf{bp\_sc\_norm}: Band-pass filtering, soft clipping, and robust per-lead normalization (z-score using median and IQR).
\end{enumerate}
Experiments read preprocessed tensors directly from disk for efficiency.

\paragraph{Band-pass filtering.}
Filtering is implemented as a linear-phase FIR band-pass using \texttt{scipy.signal.firwin}, applied with zero-phase \texttt{filtfilt} to avoid phase distortion.
Cutoffs are $0.67$--$45$\,Hz (motivated by the Zhao et al.\ ECG quality work and the defaults used in NeuroKit2/BioSPPy-style preprocessing) \todo{Add citations for Zhao et al.\ and NeuroKit2 / BioSPPy}.
The tap length is set to $\lfloor 1.5 \cdot f_s \rfloor$ with odd length enforced ($f_s=400$\,Hz after resampling); taps are reduced for short signals to fit the record length.

\paragraph{Soft clipping.}
Soft clipping is a smooth saturation transform \texttt{softclip}(x; a,b,c) applied per lead.
Lower/upper bounds $(a,b)$ are estimated from training folds only (0--3) by computing per-record $p_1/p_{99}$ and then taking the global 5th percentile of $p_1$ and 95th percentile of $p_{99}$ per lead.
The shape parameter is $c=\log(2)/\left(2/(e^2+1)\right)$ (as implemented in code), and each lead is scaled by $\max(|a|,|b|)$ to keep amplitudes comparable across leads.
\todo{find citation for this, Nabil sent me an article}

\paragraph{Robust per-lead normalization.}
The \textbf{bp\_sc\_norm} regime applies per-record, per-lead robust normalization:
\[
  x_i \leftarrow \frac{x_i - \mathrm{median}(x_i)}{\mathrm{IQR}(x_i) + \epsilon}, \quad \epsilon=10^{-6},
\]
where $\mathrm{IQR} = p_{75}-p_{25}$. No population-level statistics are used.

\section{Augmentations}
\subsection{Default augmentations}
Training uses mild augmentations to improve robustness without distorting ECG morphology:
\begin{itemize}
  \item \textbf{Amplitude scaling:} multiplicative gain jitter. \textbf{bp} uses $[0.95,1.05]$, \textbf{bp\_sc} uses $[0.975,1.025]$, and \textbf{bp\_sc\_norm} disables scaling ($[1.0,1.0]$) to preserve the intent of robust normalization.
  \item \textbf{Gaussian noise:} additive noise with $\sigma=0.005$ (per-view by default).
\end{itemize}
Stronger transforms are explicitly disabled in the main experiments: time warping (max\_time\_warp $=0$), time masking (max\_mask\_duration $=0$), channel masking (mask\_prob $=0$), and baseline-wander augmentation (wandering amplitude $=0$). This preserves interpretability and avoids injecting artifacts not supported by the target clinical setting.
\todo{Add citations for augmentation choices in ECG (e.g., mild gain/noise vs. masking/warping).}

\subsection{Physiological augmentation (rotation)}
For the \textbf{bp} regime only, we optionally apply a random rotation of the frontal electrical axis. This augmentation simulates valid physiological variation by projecting the cardiac dipole onto the limb leads at a slightly perturbed angle (implementation: approximate frontal-axis rotation in VCG space and mapping back to 12-lead ECG).
Axis rotation is enabled only in the \texttt{rot10} runs (maximum absolute rotation $10^\circ$, applied with probability 1), and kept off for \textbf{bp\_sc} and \textbf{bp\_sc\_norm} because the nonlinear preprocessing steps can break the linearity assumptions required for a physically meaningful rotation.
\todo{insert a graphic here how the signal changes per lead when rotated}

\section{Experiment Tracks}
\label{sec:track-definitions}
\subsection{Track 1: Supervised Classification}
Track~1 trains a binary classifier end-to-end. We evaluate combinations of loss functions (weighted BCE, Focal Loss, RAT), preprocessing regimes, and augmentation strategies.
Performance is primarily measured using the challenge-aligned TPR at 5\% FPR, alongside AUROC and Average Precision (AP).
We also compare against a random forest baseline trained on simple per-channel mean/std statistics to quantify the value of learning from raw waveforms.

\subsection{Track 2: Representation Learning}
Track~2 focuses on optimizing the geometric structure of the embedding space without direct classification supervision. We compare:
\begin{itemize}
  \item \textbf{SupCon (standard)}: A baseline supervised contrastive loss.
  \item \textbf{SupCon (imbalanced)}: A variant incorporating imbalance-handling logic.
  \item \textbf{Prototype-based}: An objective using learnable class prototypes.
\end{itemize}
Embedding quality is monitored using the TTC metrics defined in Section~\ref{sec:ttc-metrics} (SAD, SAA, CAC, CAD).

\subsection{Track 3: Linear Probing}
Track~3 evaluates the representations learned in Track~2 by freezing the encoder and training a simple linear classifier on top. This isolates the quality of the learned features from the end-to-end optimization capability of Track~1.

\section{Dataset Splits and Evaluation}
Given the lack of a suitable external dataset, we use a stratified 5-fold cross-validation scheme.
Splits are stratified by both patient and dataset source to prevent leakage and ensure consistent label distributions.
Folds 0--3 are used for training and validation, while fold 4 is strictly held out for final testing.

\section{Analysis Pipeline}
All post-hoc analysis is performed on the held-out test set (fold 4).
The evaluation pipeline generates:
\begin{itemize}
  \item \textbf{Clinical Metrics:} Classification performance (AUROC, TPR@5\%, AP) is calculated on the full test fold.
  \item \textbf{Embedding Diagnostics:} A fixed subset of the test fold (the "probe set") is used to compute geometric metrics (CAC, CAD) and visualize the latent space using UMAP and PCA.
  \item \textbf{Consistency Analysis:} We measure the agreement between different models using ranking correlations and intersection-over-union of the top-5\% predictions.
\end{itemize}
This ensures that all reported comparisons between tracks and preprocessing methods are based on identical test data and metric definitions.

\section{ST-DFT-LRP for XAI}
\paragraph{Implementation source.}
ST-DFT-LRP is implemented using the official DFT-LRP repository (integrated as a submodule) to ensure fidelity to the method definition.
\todo{Insert repository citation and URL.}
The method is applied to model relevance in the time domain and mapped into a short-time frequency grid using window width $H$ and shift $D$.
For interpretability and comparability, the analysis uses non-overlapping windows ($D=H$) \todo{Confirm the exact window width used (e.g., 125 samples at 400\,Hz).}

\paragraph{Qualitative per-sample composites.}
For individual ECGs, we visualize a composite view consisting of:
(i) the ST-DFT magnitude with a signed relevance overlay,
(ii) the time-domain signal with signed relevance overlay, and
(iii) the frequency spectrum with signed relevance overlay.
These panels provide a compact visual summary of where relevance concentrates in time and frequency.
\missingfigure{Composite ST-DFT-LRP visualization for representative samples from each dataset.}

\paragraph{Quantitative aggregation into 4$\times$4 grids.}
To enable comparisons across samples and groups, ST-DFT-LRP relevance is aggregated into a fixed 4$\times$4 matrix:
\begin{itemize}
  \item \textbf{Time structure (rows):} P, QRS, T, and Between-beats (assigned via NeuroKit2 delineation with midpoint-based disambiguation to ensure full coverage).
  \item \textbf{Frequency bands (columns):} 0.67--4\,Hz, 4--12\,Hz, 12--25\,Hz, 25--45\,Hz (consistent with preprocessing range).
\end{itemize}
Within each cell, relevance is summarized using a signed-balance statistic
\[
  \mathrm{balance} = \frac{\sum R}{\sum |R| + \epsilon},
\]
yielding values in $[-1,1]$ that encode directional relevance while normalizing magnitude across beats.
Beat-level matrices are normalized and averaged across beats; per-sample matrices are aggregated across leads using a relevance-mass-weighted mean.
\todo{Cite the exact delineation method (e.g., NeuroKit2 DWT) and any QC filtering applied.}
\missingfigure{Example 4$\times$4 relevance grid and group-mean comparisons.}

\paragraph{Lead diagnostics.}
For each sample, lead relevance mass is summarized to compute lead-importance statistics (per-lead mass distribution, entropy, and top-$k$ leads).
These provide a lightweight check for lead dominance effects and are used as descriptive diagnostics rather than diagnostic claims.

\chapter{Results}
\label{ch:results}

This chapter summarizes empirical findings from the evaluation pipeline described in Chapter~\ref{ch:methods}.
Unless stated otherwise, all reported metrics are computed on the held-out test fold~4.

\section{Clinical task performance (ranking under screening constraints)}
\begin{itemize}
  \item Primary endpoint: TPR@5\% (screening capacity).
  \item Supporting endpoints: pAUC@5\% FPR, AP, AUROC, and TPR@10\%.
  \item \missingfigure{Table: per-run test performance (global metrics) for all runs.}
  \item \todo{Summarize: best/worst runs and differences across preprocessing regimes and objectives.} \missingfigure{Plot: performance by preprocessing regime (small multiples or grouped bars).}
\end{itemize}

\subsection{Verified vs.\ self-reported labels (label-strength split)}
\begin{itemize}
  \item Verified labels: PTB-XL (negative) + SaMi-Trop (positive).
  \item Self-reported/mixed labels: CODE-15.
  \item \missingfigure{Table: global vs.\ verified-only vs.\ CODE-15-only metrics (TPR@5\%, pAUC@5\%, AP).}
  \item \todo{Report whether the verified-only performance drops relative to the global score and discuss possible causes (domain shift vs.\ label noise) in Chapter~\ref{ch:discussion}.} \missingfigure{Plot: verified-only vs.\ CODE-15-only performance deltas per run.}
\end{itemize}

\begin{table}[t]
  \centering
  \caption{CODE-15\% Zhao2018 quality vs.\ Chagas label (fold~4).}
  \label{tab:code15_quality_chagas}
  \begin{tabular}{lrr}
    \toprule
    \textbf{Quality}  & \textbf{Chagas=0} & \textbf{Chagas=1}                                     \\
    \midrule
    Excellent         & 241{,}976         & 3{,}796                                               \\
    Barely acceptable & 81{,}529          & 2{,}208                                               \\
    Unacceptable      & 12{,}672          & 555                                                   \\
    \midrule
    \multicolumn{3}{l}{2$\times$2 Fisher (Excellent vs.\ Not): OR=1.870, $p=3.94\times10^{-129}$} \\
    \multicolumn{3}{l}{3$\times$2 FFH Monte-Carlo: $p\approx 0.0001$}                             \\
    \bottomrule
  \end{tabular}
\end{table}
\begin{table}[t]
  \centering
  \caption{Counterfactual Zhao2018 QC distribution by applying the CODE-15\% Chagas effect to PTB-XL negatives (fold~4).}
  \label{tab:qc_counterfactual}
  \begin{tabular}{lrrr}
    \toprule
    \textbf{Quality}  & \begin{tabular}[b]{@{}r@{}}\textbf{CODE-15\% factor} \\ \textbf{(pos/neg)}\end{tabular} & \begin{tabular}[b]{@{}r@{}}\textbf{PTB-XL} \\ \textbf{neg (\%)}\end{tabular} & \begin{tabular}[b]{@{}r@{}}\textbf{PTB-XL cf} \\ \textbf{pos (\%)}\end{tabular} \\
    \midrule
    Excellent         & 0.809                                                                                   & 65.84                                                                        & 51.0                                                                            \\
    Barely acceptable & 1.329                                                                                   & 29.09                                                                        & 37.0                                                                            \\
    Unacceptable      & 2.455                                                                                   & 5.08                                                                         & 12.0                                                                            \\
    \bottomrule
  \end{tabular}
\end{table}

\section{Embedding-space diagnostics (probe set)}
\begin{itemize}
  \item Primary comparable space: encoder embeddings (\texttt{enc}) for all tracks.
  \item Metrics: GPU/CAC (collapse and global separation) and SAA (two-view consistency on the probe set).
  \item \missingfigure{Table: embedding metrics for all runs.}
  \item \todo{Summarize: which objectives/regimes show signs of collapse (GPU) and how this relates to ranking performance.} \missingfigure{Plot: GPU vs.\ TPR@5\% scatter (colored by preprocessing).}
\end{itemize}

\subsection{Probe group splits}
\begin{itemize}
  \item Probe metrics are additionally computed on:
        \begin{itemize}
          \item verified subset (PTB-XL + SaMi-Trop) vs.\ CODE-15 subset.
        \end{itemize}
  \item \todo{Summarize whether CAC/SAA differ between verified and CODE-15 subsets and how this aligns with dataset clustering.} \missingfigure{Plot: CAC/SAA split by verified vs.\ CODE-15 for each run.}
\end{itemize}

\section{Qualitative embedding visualization (PCA/UMAP hex-multipanels)}
\begin{itemize}
  \item PCA and UMAP are computed on the stored probe embeddings and visualized with hex-tiling multipanels.
  \item Panels emphasize: geometry density, Chagas enrichment, RBBB enrichment, abnormal ECG enrichment, age/$\Delta$age, and PTB-XL control flags (e.g.\ CRBBB).
  \item \missingfigure{Main figure: exemplar multipanel for selected representative runs (UMAP).}
  \item \missingfigure{Appendix: full set of multipanels for all runs (UMAP and PCA).}
  \item \todo{Describe consistent qualitative patterns (e.g.\ dataset clustering, phenotype enrichment locations) without over-interpreting.} \missingfigure{Callouts: annotated examples highlighting key regions (one per dataset).}
  \item \todo{Optional: add quantitative summaries of embedding geometry (e.g., dataset-centroid distances, within/between-dataset dispersion, or density overlap statistics) to complement the visual panels.} \missingfigure{Plot: centroid distances or dispersion ratios by dataset for UMAP/PCA.}
\end{itemize}

\section{Cross-model agreement and screening set overlap}
\begin{itemize}
  \item Global agreement: model$\times$model Spearman $\rho$ on the full-test ranking vectors.
  \item Screening agreement: top-5\% set overlap via Jaccard (IoU).
  \item Within-overlap agreement: Kendall's $\tau$ on the intersection of top-5\% sets (with intersection size reported).
  \item \missingfigure{Heatmap: Spearman rho across models.}
  \item \missingfigure{Heatmap: top-5\% IoU across models.}
  \item \todo{Summarize whether agreement clusters by preprocessing regime, loss/objective, or track.}
\end{itemize}

\section{Consensus sets for sample selection (preparation for DFT-LRP)}
\begin{itemize}
  \item Models define top-5\% sets $T_m$; per-sample consensus score is $c_i=\sum_m \mathbf{1}(i\in T_m)$.
  \item Candidate tables enforce patient-unique sampling and retain QC metadata for manual selection.
  \item \missingfigure{Table: examples of high-consensus, low-consensus, and non-consensus candidates (with dataset and QC).}
  \item \todo{State final thresholds $\alpha,\beta$ used for consensus-high/low (if used) and the rationale.} \missingfigure{Histogram: distribution of consensus scores c\_i (and normalized consensus).}
\end{itemize}

\subsection{Robust rank aggregation (RRA) by dataset}
RRA \parencite{rra_kolde_2012}is applied to the full-test rankings to identify samples that are \emph{robustly} ranked high or low across models using the python package \texttt{PyFLAGR}\parencite{pyflagr_2023} (version 1.0.21).
The total test set size was $N=73{,}236$, with $7{,}411$ robust positives and $11{,}456$ robust negatives.
Figure~\ref{fig:rra_dataset_violin} summarizes the rank distributions for the robust positive and robust negative sets by dataset.
RRA provides an ordering \emph{within} each robust set; samples not included in the robust lists are not ranked by RRA.

\begin{wrapfigure}{r}{0.7\textwidth}
  \centering
  \includegraphics[width=\linewidth]{../analysis/embeddings_probe/ranking_agreement/test/plots/rra_ridgeline.png}
  \caption{RRA rank distributions by dataset and robust set. Each ridge shows the distribution of RRA ranks for robust positives or robust negatives within a dataset (lower rank = more robust).}
  \label{fig:rra_dataset_violin}
\end{wrapfigure}

\begin{table}[t]
  \centering
  \caption{Kruskal--Wallis tests for RRA rank distributions across datasets.}
  \label{tab:rra_kw}
  \begin{tabular}{lrr}
    \toprule
    \textbf{Subset}  & \textbf{H statistic} & \textbf{p-value}     \\
    \midrule
    Robust positives & 147.2232             & $1.07\times10^{-32}$ \\
    Robust negatives & 90.1182              & $2.70\times10^{-20}$ \\
    \bottomrule
  \end{tabular}
\end{table}

\section{Summary of key results}
\begin{itemize}
  \item \todo{One paragraph: main takeaways for (i) ranking performance, (ii) embedding health, (iii) dataset/label effects, (iv) agreement/consensus.} \missingfigure{One consolidated summary figure (optional): performance vs.\ embedding metrics.}
\end{itemize}

\chapter{Discussion}
\label{ch:discussion}

This thesis compares supervised classification (Track~1), representation learning (Track~2), and downstream probing of learned representations (Track~3).
While the same backbone architecture is used throughout for consistency, the tracks answer different questions and should be interpreted accordingly.

\paragraph{Linear probe vs.\ end-to-end classification}
Track~3 is implemented as a linear evaluation protocol: an encoder pretrained in Track~2 is reused and (by default) frozen, and only a lightweight classifier head is trained to predict the Chagas label.
This probes whether label information is readily accessible in the learned representation without allowing the encoder to co-adapt to the downstream objective.
In contrast, Track~1 trains encoder and classifier jointly and therefore reflects the best achievable supervised classifier under the chosen objective and preprocessing.
Consequently, Track~1 and Track~3 are not an apples-to-apples comparison of classifier performance; Track~3 is primarily used to connect improvements in embedding-space diagnostics to downstream separability.

\paragraph{Clinical label noise and short-window ECGs.}
The labels in this thesis are based on seropositivity, yet clinical staging defines the indeterminate form (Stage A) as seropositive with a normal ECG and early cardiomyopathy (Stage B1) as the first appearance of ECG abnormalities in a seropositive patient \parencite{chagas_progression_2022,lancet_clinical_chagas_features_2024}.
This creates biological label noise for ECG-based classification because a large fraction of positives may be ECG-normal, while the model's signal is driven by electrophysiological manifestations such as RBBB/LAFB and ventricular ectopy that are typical of CCC \parencite{chagas_ecg_abnormalities_2018,lancet_clinical_chagas_features_2024}.
In addition, the largest dataset (CODE-15\%) uses self-reported Chagas labels, which introduces extra uncertainty beyond staging effects, whereas SaMi-Trop provides clinically confirmed positives \parencite{ribeiro_code-15_2021,ribeiro_sami-trop_2021}.
Conversely, PTB-XL provides comparatively clean negatives because it was recorded in Germany (1989--1996), a non-endemic setting, so Chagas labels are treated as negatives on epidemiological grounds \parencite{wagner_ptb_xl}.
Given short (~6.25~s) windows, the models are inherently biased toward beat morphology and high-burden ectopy rather than low-frequency paroxysmal events, which helps explain why performance is stronger in cohorts enriched for cardiomyopathy and why ranking-based screening is the most clinically plausible objective.

\paragraph{Augmentation strength as a limiting factor.}
The augmentation suite used in this thesis was intentionally mild (small gain jitter and low-variance Gaussian noise), with stronger perturbations such as time warping, time masking, and channel masking disabled to preserve clinical interpretability.
While this choice stabilizes supervised training and supports direct waveform interpretation, it may be suboptimal for contrastive representation learning: Track~2 objectives generally benefit from stronger or more diverse view generation to encourage invariances.
Therefore, some of the observed limits in Track~2 (and downstream Track~3) performance could reflect underpowered augmentations rather than intrinsic limitations of the objectives.

\paragraph{Lead construction and rotation fidelity.}
All experiments operate on the 12-lead signals as provided by the source datasets; no new leads are re-derived from raw electrode potentials, and the only signal modifications are the defined preprocessing and augmentations.
Consequently, the axis-rotation augmentation relies on a fixed linear lead model applied to potentially heterogeneous acquisition pipelines and gain calibrations.
This means the rotation parameters are approximate and may not correspond to a precise physical change in cardiac orientation for every record, which is another reason the rotation should be interpreted as a controlled robustness probe rather than a clinically exact transformation.

\paragraph{Lead selection in ST-DFT-LRP.}
Beat-level ST-DFT-LRP aggregation in this thesis is restricted to Lead~II, which is a standard rhythm lead and provides a clear P--QRS--T morphology.
This choice improves feasibility and interpretability but is not optimal for all Chagas-related abnormalities: several structural and conduction manifestations (e.g., bundle-branch patterns or regional ST--T changes) are known to be more prominent in precordial or alternative limb leads.
Therefore, the lead-II focus should be interpreted as a rhythm-centric view of model attention rather than a comprehensive 12-lead characterization of Chagas cardiomyopathy.

\paragraph{RBBB annotations across datasets.}
For CODE15, RBBB flags are provided directly in the dataset metadata; for PTB-XL, the closest available analogue used here is \texttt{ptb\_crbbb} derived from \texttt{scp\_codes}.
This mapping allows coarse, cross-dataset comparisons of bundle-branch related patterns, but should be interpreted cautiously because complete RBBB in PTB-XL is not a perfect substitute for the broader RBBB labels in CODE15.

\paragraph{Architecture-specificity and generalization.}
All results in this thesis are conditional on a fixed ResNet-18 backbone, chosen to isolate the effects of preprocessing and objective functions rather than to optimize architecture.
This implies that the reported gains and failure modes are most directly applicable to CNN-style inductive biases that emphasize local morphology and translation invariance.
Architectures with different temporal aggregation (e.g., attention-based models or dilated TCNs) may emphasize rhythm, interval structure, or long-range dependencies differently, which can change the sensitivity to augmentation, normalization, and dataset bias. \todo{List 1--2 alternative architecture families relevant to ECGs used in your literature scan.}
In particular, the axis-rotation augmentation may interact differently with models that fuse leads late or model inter-lead correlations explicitly. \todo{Note any lead-fusion strategy differences observed in prior ECG architectures.}
Future work should therefore validate the key findings (e.g., loss-function effects, preprocessing regimes, rotation robustness) on at least one non-CNN baseline to assess architectural transferability. \todo{Specify a concrete baseline to test (e.g., 1D Transformer, InceptionTime, or a dilated TCN).}
Future work should revisit augmentation strength in Track~2 specifically, exploring moderate temporal perturbations or structured masking while validating that ECG morphology remains interpretable and physiologically plausible.

\todo{beobachtungen nur für diese resnet18 architektur?}


\appendix
\chapter{Appendix}

Put supplementary material here.



\clearpage
\printbibliography

\end{document}
